\title{QLoRA: Efficient Finetuning of Quantized LLMs}

\begin{abstract}
We introduce \textsc{QLoRA}, a novel and efficient finetuning method designed to drastically lower memory demands, enabling the finetuning of a 65B parameter language model on a single 48GB GPU without any loss in the effectiveness typically achieved by standard 16-bit finetuning. In \textsc{QLoRA}, gradients are propagated through a fixed, 4-bit quantized version of a pretrained language model and into Low Rank Adapters (LoRA). Our leading model suite, termed \textbf{Guanaco}, achieves superior results on the Vicuna benchmark, surpassing all previously available open models and attaining 99.3\% of ChatGPT’s performance after just 24 hours of single-GPU finetuning. \textsc{QLoRA} brings several technical advancements that allow for substantial memory savings without diminishing model capability: (a) 4-bit NormalFloat (NF4), a newly proposed data format theoretically optimized for weights with normal distributions; (b) Double Quantization, which compresses memory further by quantizing the quantization coefficients themselves; and (c) Paged Optimizers, which help control transient increases in memory consumption. Utilizing \textsc{QLoRA}, we have finetuned over 1,000 models and systematically studied instruction-following and conversational agent performance across 8 instruction datasets, multiple architectures (LLaMA, T5), and various large model sizes that are ordinarily inaccessible with traditional finetuning approaches (such as 33B and 65B parameter models). Our findings demonstrate that employing QLoRA with compact, high-quality datasets produces leading results—even with smaller models compared to prior state-of-the-art methods. We provide comprehensive performance analyses using both human raters and GPT-4, highlighting that GPT-4-based evaluations are a practical, low-cost proxy for human judgment. Additionally, our examination reveals the unreliability of existing chatbot benchmarks in genuinely assessing model capability. Through a “lemon-picked” case study, we identify and illustrate specific shortcomings where \textbf{Guanaco} does not match ChatGPT. We are making available all models, code, and CUDA kernels required for 4-bit training.\footnote{\url{https://github.com/artidoro/qlora} and \url{https://github.com/TimDettmers/bitsandbytes}}
\end{abstract}

\section{Introduction}

Adapting large language models (LLMs) through finetuning remains an exceptionally powerful strategy to enhance model capabilities \citep{min2021metaicl, wei2021finetuned, ouyang2022training, wang2022super, wang2022self, liu2022few} and to tailor or eliminate specific model behaviors \citep{ouyang2022training,askell2021general,bai2022training}. Nevertheless, finetuning extensive models comes with immense computational demands; for example, conventional 16-bit finetuning for the LLaMA model with 65B parameters~\citep{touvron2023llama} consumes over 780 GB of GPU memory. While recent progress in quantization has led to techniques capable of reducing LLM memory consumption for inference \citep{dettmers2022llm, dettmers2022case, frantar2022gptq, xiao2022smoothquant}, such approaches are limited to inference and are not robust during the training phase \citep{wortsman2023stable}.

In this work, we show for the first time that a quantized 4-bit model can be finetuned without any sacrifice in accuracy. Our approach, \textsc{QLoRA}, introduces a novel quantization mechanism that transforms a pretrained model to 4-bit precision with high fidelity, and then incorporates a lightweight, trainable set of Low-rank Adapter weights \citep{hu2021lora} 

which are optimized by propagating gradients through the quantized layers.

With \textsc{QLoRA}, the finetuning memory requirement for a 65B parameter model is reduced dramatically—from over 780GB to less than 48GB—while preserving both speed and accuracy relative to standard 16-bit finetuning. This development significantly broadens the feasibility of LLM finetuning, making it possible to finetune the largest open-source models available today on a single GPU. Leveraging \textsc{QLoRA}, we introduce the \textbf{Guanaco} series of models; notably, our second-best model attains 97.8\% of ChatGPT’s performance on the Vicuna benchmark~\citep{vicuna2023}, requiring under 12 hours of training on an off-the-shelf consumer GPU. Training on a single high-end professional GPU for 24 hours, our largest model achieves 99.3\%—essentially matching ChatGPT on the same benchmark. Additionally, our smallest \textbf{Guanaco} model, with 7B parameters, delivers strong results, needing only 5 GB of memory and outperforming the 26 GB Alpaca model by more than 20 percentage points on the Vicuna evaluation (see Table~\ref{tbl:vicuna_eval}).

\textsc{QLoRA} incorporates several key advancements that collectively reduce memory requirements while maintaining high performance: (1) \textbf{4-bit NormalFloat}, an optimally information-preserving quantization format tailored for data following a normal distribution, which empirically surpasses the commonly used 4-bit Integers and 4-bit Floats. (2) \textbf{Double Quantization}, a technique that further compresses the quantization parameters themselves, achieving an average reduction of approximately 0.37 bits per parameter—amounting to roughly 3 GB of savings for a 65B parameter model. (3) \textbf{Paged Optimizers}, which exploit NVIDIA’s unified memory infrastructure to mitigate the memory spikes caused by gradient checkpointing during the processing of long-sequence mini-batches. These innovations are combined within a more refined LoRA framework, characterized by the integration of adapters at every layer of the network, thereby eliminating most of the accuracy compromises observed in previous approaches.

Owing to the memory efficiency of \textsc{QLoRA}, we are able to undertake comprehensive experiments examining instruction fine-tuning and chatbot effectiveness across model scales that would not be feasible with conventional fine-tuning methods because of their memory demands. Specifically, we train over 1,000 models spanning multiple instruction-tuning corpora, diverse architectures, and sizes ranging from 80M to 65B parameters. Beyond demonstrating that \textsc{QLoRA} attains the same level of performance as 16-bit finetuning (\S\ref{sec:comp_to_std}) and facilitating the training of a state-of-the-art chatbot, \textbf{Guanaco} (\S\ref{sec:pushing}), we also investigate patterns within the resulting models. Our findings highlight that dataset quality outweighs dataset size: for instance, a compact 9k-example set (OASST1) yielded better chatbot results than a much larger, subsampled 450k-example set (FLAN v2), even when both targeted generalization to instruction following. Additionally, our results reveal that strong performance on the Massive Multitask Language Understanding (MMLU) benchmark does not necessarily translate to superior results on the Vicuna chatbot benchmark, and vice versa—emphasizing that the appropriateness of the dataset for the intended task has a greater impact than its size.

In addition, we conduct a thorough investigation of chatbot capabilities, incorporating both assessments by human judges and GPT-4 as evaluation mechanisms. For our methodology, we utilize a tournament-based framework in which models engage in head-to-head competitions to generate optimal responses to specified prompts. Each match’s outcome is determined by either human annotators or GPT-4. We synthesize tournament outcomes into Elo scores~\citep{elo1967proposed, elo1978rating}, which establish the relative ranking of chatbot models. Our analysis shows a high degree of concordance between GPT-4 and human judgments with respect to model rankings, though notable points of substantial disagreement emerge. These findings suggest that, while automated evaluation offers a cost-efficient alternative to manual annotation, it introduces certain ambiguities.

To supplement the quantitative benchmarking, we present an in-depth qualitative examination of the \textbf{Guanaco} models. This qualitative assessment sheds light on instances of both strength and weakness that are not reflected in the numerical evaluation.

To promote continued research, we are releasing the model outputs alongside both human and GPT-4 annotation labels. The entirety of our codebase and CUDA kernel implementations are open-sourced, and we have incorporated our procedures into the Hugging Face transformers library \citep{wolf2019huggingface}, ensuring widespread usability. Furthermore, we provide a set of adapters compatible with 7/13/33/65B scale models, trained across eight distinct instruction-following datasets—yielding 32 openly available, fine-tuned models.

\section{Background}
\paragraph{Block-wise k-bit Quantization} 
Quantization refers to transforming data from a higher-information representation to one that conveys less information. Practically, this process typically involves converting data from a format with a large bit width (such as 32-bit floating point) into one with fewer bits (like 8-bit integers). To fully utilize the expressiveness of the lower-bit datatype, the original input—often arranged as a tensor—is rescaled into the smaller range, typically by normalizing each value according to the largest absolute value in the input. For instance, to quantize a tensor of 32-bit floats (FP32) into a signed 8-bit integer (Int8) tensor with permissible values in $[-127, 127]$:

\begin{equation}
    \mathbf{X}^{ \text{Int8}} = \text{round}\left( \frac{127}{{\text{absmax}}(\mathbf{X}^{{\text{FP32}}})}\mathbf{X}^{{\text{FP32}}} \right) = \text{round}(c^{\text{FP32}}\cdot \mathbf{X}^{{\text{FP32}}}),
\end{equation}

Here, $c$ serves as the {\em quantization constant} or {\em scale}. Reversing the quantization—dequantization—retrieves the original values:

\begin{equation}
    \text{dequant}(c^{\text{FP32}}, \mathbf{X}^{\text{Int8}}) = \frac{\mathbf{X}^{\text{Int8}}}{c^{\text{FP32}}} = \mathbf{X}^{\text{FP32}}
\end{equation}

A limitation of this strategy emerges when the input tensor contains values of exceptionally high magnitude (outliers), which results in inefficient use of quantization bins: certain bit patterns correspond to bins that end up nearly empty or unused altogether. To address the challenge posed by outliers, a widely adopted technique involves dividing the input tensor into several blocks and quantizing each block independently using its own quantization parameter $c$. More precisely, the input tensor $\mathbf{X}\in \mathbb{R}^{b\times h}$ is first reshaped into a contiguous sequence and then partitioned into $n$ blocks of size $B$, such that the total number of blocks is $n = ({b\times h})/{B}$. Quantization is subsequently performed separately on each block, following Equation~1, resulting in a quantized tensor accompanied by a unique quantization constant $c_i$ for each block.

\paragraph{Low-rank Adapters} LoRA finetuning \citep{hu2021lora} introduces memory savings by employing a compact set of trainable parameters, known as adapters, while keeping the principal model weights fixed during training. During stochastic gradient descent, only the adapters receive parameter updates; gradients propagate through the frozen pretrained weights directly to the adapters, which are trained to minimize the objective. LoRA implements a modification to standard linear projections by incorporating an additional low-rank, factorized projection. For a standard linear operation ${\mathbf{X} \mathbf{W} = \mathbf{Y}}$ where $\mathbf{X}\in  \mathbb{R}^{b\times h}$ and $\mathbf{W}\in  \mathbb{R}^{h\times o}$, LoRA applies:

\begin{equation}
    \mathbf{Y} = \mathbf{X}\mathbf{W} + s\mathbf{X}\mathbf{L}_1\mathbf{L}_2,
\end{equation}

where the adapter weights $\mathbf{L}_1\in  \mathbb{R}^{h\times r}$ and $\mathbf{L}_2\in  \mathbb{R}^{r\times o}$ capture low-rank information, and $s$ is a scaling factor.

\paragraph{Memory Requirement of Parameter-Efficient Finetuning} The memory demand associated with LoRA during training warrants careful consideration, particularly regarding both the quantity and dimensionality of adapters deployed. Owing to LoRA’s exceptionally small memory overhead, it becomes feasible to incorporate a greater number of adapters to boost model performance while only marginally impacting the aggregate memory usage. Although LoRA belongs to the broader class of Parameter Efficient Finetuning (PEFT) approaches, the dominant memory usage in LLM finetuning arises from the storage of activation gradients, not the LoRA parameters themselves. For instance, consider a 7B LLaMA model trained on FLAN v2 using a batch size of 1 and LoRA weights comprising the standard 0.2\% of total model weights \citep{hu2021lora,liu2022few}. Here, the memory needed for LoRA input gradients totals 567 MB, compared to a mere 26 MB occupied by LoRA parameters. Employing gradient checkpointing~\citep{chen2016training} reduces the memory per sequence for input gradients to about 18 MB, yet this still exceeds the storage requirement of all LoRA parameters combined. In context, the 4-bit quantized base model uses 5,048 MB of memory. These findings emphasize not only the value of gradient checkpointing but also the limited memory savings gained from reducing the number of LoRA parameters further. Consequently, increasing the number of adapters remains possible without significantly augmenting the memory needed for training (see Appendix~\ref{app:memory} for a comprehensive analysis). As elaborated in subsequent sections, this property is vital for achieving performance equivalent to models trained with full 16-bit precision.

\section{\textsc{QLoRA} Finetuning}

\textsc{QLoRA} facilitates precise 4-bit finetuning by leveraging two core innovations we introduce: 4-bit NormalFloat (NF4) quantization and Double Quantization. To further address memory management challenges—particularly out-of-memory (OOM) incidents during gradient checkpointing, which have traditionally limited single-machine finetuning for large-scale models—we present Paged Optimizers as an effective mitigation.

\textsc{QLoRA} operates with two principal data types: a storage data type, generally set to 4-bit precision, and a separate computation data type, commonly BFloat16. Practically, this setup entails that whenever a weight tensor in \textsc{QLoRA} is accessed, it undergoes dequantization into the BFloat16 format prior to engaging in 16-bit matrix multiplication operations.

We proceed by detailing the primary constituents of \textsc{QLoRA}, culminating in a formal articulation of its methodology.

\paragraph{4-bit NormalFloat Quantization} The NormalFloat (NF) format extends Quantile Quantization \citep{dettmers2022optimizers}, a data type underpinned by information-theoretical optimality that assigns an identical number of tensor values to each quantization bin. Quantile quantization accomplishes this by calculating the quantiles of an input tensor using its empirical cumulative distribution function.

Despite its theoretical appeal, quantile quantization is hindered by the computational cost of accurately estimating quantiles. As a remedy, swift quantile approximation techniques such as SRAM quantiles~\citep{dettmers2022optimizers} are employed. Nevertheless, these approximations often result in substantial quantization errors, particularly for outlier values that may hold significant importance.

Such computational overheads and approximation-induced inaccuracies can be circumvented when the input tensor follows a distribution consistent up to a quantization constant; in these scenarios, the quantiles across tensors remain invariant, making precise quantile computation both viable and efficient.

Pretrained neural network weights are typically distributed according to a zero-mean normal distribution characterized by a standard deviation $\sigma$ (refer to Appendix~\ref{app:normality}). To standardize these weights across different scenarios, we can rescale $\sigma$ so that the weights' distribution aligns precisely with the desired numerical range. In this work, we designate $[-1, 1]$ as the operative interval for our data type. Consequently, it is necessary to map both the data type’s quantiles and the distribution of neural network weights into this unified interval.

The most information-efficient data type for representing zero-mean normal distributions with variable $\sigma$ over the interval $[-1, 1]$ can be constructed as follows: (1) first, compute the $2^k+1$ quantiles of the standard normal distribution $N(0, 1)$ to derive a quantile-based $k$-bit quantization scheme, (2) scale the resulting data type’s quantile values so they fit exactly within the $[-1, 1]$ range, and (3) quantize the given weight tensor by rescaling it to $[-1, 1]$ via an absolute maximum normalization.

After aligning the weight tensor’s range with that of the data type, standard quantization procedures can be applied. The rescaling in step (3) effectively adjusts the standard deviation of the weight tensor to equal that of the $k$-bit data type. Formally, the $2^k$ representative values $q_i$ for the data type are estimated by

\begin{equation}
q_i  = \frac{1}{2}\left( Q_X\left(\frac{i}{2^k+1}\right) + Q_X\left(\frac{i+1}{2^k+1}\right)\right),
\end{equation}

where $Q_X(\cdot)$ denotes the quantile function for the standard normal distribution $N(0,1)$. A limitation of symmetric $k$-bit quantization is its inability to exactly represent zero, which is crucial for lossless quantization of padding and other elements inherently assigned a value of zero. To guarantee a discrete zeropoint at $0$ and utilize the full range of $2^k$ bit patterns in a $k$-bit format, we instead construct an asymmetric data type. This is done by separately calculating the quantiles $q_i$ for two regions: allocating $2^{k-1}$ bins for negative values and $2^{k-1}+1$ bins for positive values. These sets of $q_i$ are then merged, and the duplicated zero entry, arising in both partitions, is removed. We refer to this construction, which ensures an equal expected population across all quantization intervals, as the \emph{k-bit NormalFloat} (NFk) data type. NFk is information-theoretically optimal for quantizing zero-centered normal distributions. Details on the precise quantized values for this data type are provided in Appendix~\ref{app:nf4}.

\paragraph{Double Quantization{}}
We propose \emph{Double Quantization} (DQ), a technique that quantizes the quantization constants themselves to further conserve memory. Accurate 4-bit quantization requires small block sizes~\citep{dettmers2022case}, yet such small blocks incur nontrivial memory overhead. As an illustration, storing quantization constants in 32 bits with a block size of 64 for $\mathbf{W}$ results in an average cost of $32/64 = 0.5$ bits per parameter. Double Quantization addresses this overhead by compressing the quantization constants further.

Concretely, Double Quantization involves taking the initial quantization constants $c_2^{\text{FP32}}$ as the target for a second quantization stage. The outcome is a quantized set of constants $c_2^{\text{FP8}}$, together with new quantization constants $c_1^{\text{FP32}}$ at the second level. We employ 8-bit floating-point representations with a block size of 256 in this stage, as prior findings from \citet{dettmers2022case} indicate that this granularity maintains performance. Since $c_2^{\text{FP32}}$ are strictly positive, we mean-center $c_2$ prior to quantization, enabling the use of symmetric quantization methods. As a result, for a block size of 64, this two-level quantization reduces the memory use per parameter from $32/64 = 0.5$ bits to $8/64 + 32/(64\cdot256) = 0.127$ bits, amounting to a savings of 0.373 bits per parameter.

\paragraph{Paged Optimizers} leverage NVIDIA's unified memory~\footnote{\tiny \url{https://docs.nvidia.com/cuda/cuda-c-programming-guide}}, a mechanism that transparently manages memory transfers between the CPU and GPU at the level of individual memory pages, thereby preventing GPU memory errors when the device memory is exhausted. This feature operates analogously to typical memory paging, where data is swapped between CPU RAM and storage. In our approach, we utilize this unified memory capability to allocate optimizer state variables as paged memory. Consequently, when GPU memory limits are reached, the optimizer state is automatically moved to CPU RAM and subsequently reloaded into GPU memory as required during the optimizer's update operations.

\paragraph{\textsc{QLoRA}.} Building upon the aforementioned components, we construct \textsc{QLoRA} for an individual linear layer in a quantized backbone model that employs a single LoRA adapter, described by the following:

\begin{equation}
    \mathbf{Y}^{\text{BF16}} =  \mathbf{X}^{\text{BF16}} \text{doubleDequant}(c_1^{\text{FP32}}, c_2^{\text{k-bit}}, \mathbf{W}^{\text{NF4}}) + \mathbf{X}^{\text{BF16}}\mathbf{L}^{\text{BF16}}_1\mathbf{L}^{\text{BF16}}_2,
\end{equation}

The function doubleDequant$(\cdot)$ is further defined as:

\begin{equation}
    \text{doubleDequant}(c_1^{\text{FP32}}, c_2^{\text{k-bit}}, \mathbf{W}^{\text{k-bit}}) = \text{dequant}(\text{dequant}(c_1^{\text{FP32}}, c_2^{\text{k-bit}}), \mathbf{W}^{\text{4bit}}) = \mathbf{W}^{\text{BF16}},
\end{equation}

We employ NF4 quantization for $\mathbf{W}$ and adopt FP8 format for $c_2$.

A blocksize of 64 is selected for $\mathbf{W}$ to achieve finer quantization accuracy, while a blocksize of 256 is used for $c_2$ in order to minimize memory overhead.

When updating parameters, only gradients with respect to the adapter weights, $\frac{\partial E}{\partial \mathbf{L}_i}$, are computed; gradients for the 4-bit weights, $\frac{\partial E}{\partial \mathbf{W}}$, are not required. Nevertheless, deriving $\frac{\partial E}{\partial \mathbf{L}_i}$ involves the evaluation of $\frac{\partial \mathbf{X}}{\partial \mathbf{W}}$. This requires dequantizing the stored $\mathbf{W}^{\text{NF4}}$ back into the BFloat16 representation $\mathbf{W}^{\text{BF16}}$—as illustrated by equation~(5)—to obtain $\frac{\partial \mathbf{X}}{\partial \mathbf{W}}$ using BFloat16 precision.

In summary, \textsc{QLoRA} utilizes a single storage data type—most often 4-bit NormalFloat—and a distinct computation data type, namely 16-bit BrainFloat. During training, weights are dequantized from the storage format to the computation format to carry out both the forward and backward passes; however, only the LoRA parameters—which are kept in 16-bit BrainFloat—are subject to gradient computations.

\section{QLoRA vs. Standard Finetuning}
\label{sec:comp_to_std}

Having outlined the inner workings of QLoRA and its substantial potential for lowering the memory footprint during model finetuning, a critical issue is whether QLoRA achieves comparable performance to that of conventional full-model finetuning. Moreover, it is important to assess the influence of QLoRA’s individual elements, particularly the effect of NormalFloat4 compared to traditional Float4. The following subsections present the experiments devised to address these topics.

\paragraph{Experimental setup.} The study involves three neural architectures—encoder, encoder-decoder, and decoder-only models—and examines the performance of QLoRA relative to 16-bit adapter finetuning as well as full-model finetuning, with experiments carried out on models up to 3B parameters. Evaluations span a variety of benchmarks: GLUE \citep{wang2018glue} using RoBERTa-large \citep{liu2019roberta}; Super-NaturalInstructions (TKInstruct) \citep{wang2022super} with T5 \citep{t5}; and 5-shot MMLU \citep{hendrycksmeasuring} after adapting LLaMA to Flan v2 \citep{longpre2023flan} and Alpaca \citep{alpaca}. In order to directly investigate the benefits of NF4 over alternative 4-bit data formats, the framework proposed by \citet{dettmers2022case} is employed to evaluate post-quantization zero-shot accuracy and perplexity for a range of models (OPT~\citep{zhang2022opt}, LLaMA~\citep{touvron2023llama}, BLOOM~\citep{scao2022bloom}, Pythia~\citep{biderman2023pythia}) covering model scales from 125 million to 13 billion parameters.

Additional details for each experimental context are discussed in the results sections to enhance clarity, with comprehensive setup information available in Appendix~\ref{app:exp-setup}.

While paged optimizers are indispensable for enabling 33B/65B \textsc{QLoRA} finetuning on a single 24/48GB GPU, we have not systematically reported precise metrics for the impact of paged optimizers. This is due to the infrequent nature of paging events, which typically only occur with very long input sequences. Nevertheless, we include a runtime analysis involving 65B-parameter models on 48GB GPUs, demonstrating that for a batch size of 16, paged optimizers yield training throughput on par with that of standard optimizers. Identifying and quantifying scenarios in which paging leads to notable slowdowns remains an open area for future research.

\paragraph{Default LoRA hyperparameters do not match 16-bit performance} 

Adopting the commonly recommended strategy of restricting LoRA adaptation to only the query and value matrices in the attention layers \citep{hu2021lora} falls short in reproducing the performance of full finetuning, particularly for larger base models. Evidence from Figure~\ref{fig:hyper}, which presents results from LLaMA 7B finetuning on Alpaca, suggests that the determining factor among LoRA hyperparameters is the overall count of LoRA adapters. Achieving parity with full finetuning requires applying LoRA to all linear layers within the transformer blocks. Other hyperparameters, like the projection rank $r$, were not found to influence performance substantially (see Appendix \ref{app:exp-setup} for more information).

In a similar vein, our analysis reveals that the default hyperparameters used for fully finetuned baselines are insufficiently optimized. To address this, we conduct an extensive hyperparameter search, exploring learning rates ranging from 1e-6 to 5e-5 and batch sizes between 8 and 128, with the aim of establishing reliable baselines. Figure~\ref{fig:hyper} presents the outcomes of finetuning the 7B LLaMA model on the Alpaca dataset.

\paragraph{4-bit NormalFloat Outperforms 4-bit Floating Point}

Although the 4-bit NormalFloat (NF4) representation is optimal from an information-theoretic perspective, it remains an open question whether this translates into practical improvements. We adopt the experimental protocol outlined by \citet{dettmers2022case}, evaluating quantized LLMs (including OPT \citep{zhang2022opt}, BLOOM \cite{scao2022bloom}, Pythia \cite{biderman2023pythia}, and LLaMA) at varying scales (125M–65B parameters) and employing a range of quantization types for language modeling and several zero-shot evaluation tasks. According to Figure~\ref{fig:nf4} and Table~\ref{tbl:nf4_ppl}, NF4 leads to notable improvements compared to FP4 and Int4, and employing double quantization can further reduce memory usage without sacrificing accuracy.

\paragraph{k-bit \textsc{QLoRA} Achieves Parity with 16-bit Full and LoRA Finetuning} 

Recent studies indicate that while 4-bit quantization is feasible for \emph{inference}, it typically results in some performance loss relative to 16-bit models \citep{dettmers2022case,frantar2022gptq}. This raises the critical question of whether this performance deficit can be addressed through 4-bit adapter finetuning. To investigate this, we conduct experiments in two configurations.

The first scenario involves a direct comparison to full 16-bit finetuning for RoBERTA and T5 architectures, with parameter sizes from 125M to 3B, on both the GLUE benchmark and the Super-NaturalInstructions dataset. The findings, detailed in Table~\ref{tab:nf4vsbf16}, demonstrate that adapter-based finetuning using 16-bit, 8-bit, or 4-bit precision yields results indistinguishable from those of full 16-bit finetuning. This provides evidence that performance losses due to coarse quantization can be fully mitigated by applying adapter-based finetuning post-quantization.

In our second experimental configuration, fully finetuning models with 11B parameters or more necessitates multiple high-memory GPU servers. Therefore, we further investigate whether 4-bit \textsc{QLoRA} can achieve parity with 16-bit LoRA within the 7B to 65B parameter range. Specifically, we finetune LLaMA models ranging from 7B to 65B parameters on two instruction-tuned datasets, Alpaca and FLAN v2, subsequently evaluating performance on the MMLU benchmark using 5-shot accuracy. The outcomes, detailed in Table~\ref{tbl:llama-datatype}, demonstrate that NF4 with double quantization is able to fully recapture the MMLU results of 16-bit LoRA. Furthermore, our analysis reveals that \textsc{QLoRA} using FP4 falls approximately 1 percentage point short of the baseline set by 16-bit brain float LoRA. These observations reinforce two primary conclusions: (1) \textsc{QLoRA} employing NF4 attains results equivalent to both 16-bit full and LoRA finetuning; and (2) among quantization options, NF4 surpasses FP4 in terms of precision.

\paragraph{Summary} Our experimental findings consistently indicate that applying 4-bit \textsc{QLoRA} with the NF4 data type produces finetuning results on par with 16-bit full and LoRA approaches on academic benchmarks evaluated under established methodologies. The results further illustrate that NF4 outperforms FP4, and that double quantization does not result in a drop in model performance. Together, these points offer substantial support for the reliability of 4-bit \textsc{QLoRA} in matching the effectiveness of 16-bit finetuning methods.

Consistent with previous research on quantization \citep{dettmers2022case}, our analyses of MMLU and Elo scores show that, for a fixed allocation of finetuning and inference resources, it is advantageous to expand the parameter count of the base model while reducing precision. This finding underscores the practical efficiency benefits provided by \textsc{QLoRA}. As we observed no significant loss in performance relative to full finetuning when employing 4-bit quantized models, a crucial question that remains is to pinpoint the precise boundary of the performance-precision trade-off in QLoRA tuning, which we identify as a direction for future investigation.

We next explore instruction tuning at scales that would be unattainable with complete 16-bit finetuning, particularly given the constraints of typical academic computational resources.

\section{Pushing the Chatbot State-of-the-art with QLoRA}
\label{sec:pushing}
After demonstrating that 4-bit \textsc{QLoRA} consistently achieves parity with 16-bit finetuning across various model sizes, tasks, and datasets, we proceed with a comprehensive analysis of instruction finetuning applied to the largest open-source language models accessible for academic study. To evaluate how instruction finetuning influences these models, we utilize a demanding Natural Language Understanding benchmark (MMLU) and introduce novel methodologies for systematically measuring chatbot performance in realistic scenarios.

\subsection{Experimental setup} Here, we present a summary of the experimental methodology; further details are provided in Appendix \ref{app:chatbot}.

\paragraph{Data} Since, to our knowledge, there lacks an exhaustive comparative analysis of the latest instruction-following datasets, we curate a selection of eight contemporary datasets for our experiments. Our choices span datasets derived from crowd-sourcing (OASST1~\citep{kopf2023openassistant}, HH-RLHF~\citep{bai2022training}), distillation processes involving instruction-tuned models (Alpaca~\citep{alpaca}, self-instruct~\citep{wang2022self}, unnatural-instructions~\citep{honovich2022unnatural}), comprehensive corpora aggregations (FLAN v2~\citep{chung2022scaling}), and hybrid approaches (Chip2~\citep{laion2023}, Longform~\citep{koksal2023longform}). Collectively, these datasets encompass a diverse array of languages, dataset sizes, and licensing conditions.

\paragraph{Training Setup} To minimize confounding factors related to varying training objectives, all QLoRA finetuning is conducted using cross-entropy loss (supervised learning), explicitly excluding reinforcement learning techniques—even for data containing human evaluation of responses. For datasets where instructions and responses are readily distinguishable, the finetuning is applied only to the response portion (see Appendix \ref{app:chatbot} for ablation studies). When using OASST1 and HH-RLHF, which provide multiple candidate responses, we choose the highest-ranked response at each conversational node and finetune on the entire corresponding conversation thread, including the instructions. Throughout, NF4 \textsc{QLoRA} is employed with double quantization and paged optimizers to avoid memory overhead during gradient checkpointing. For hyperparameter optimization, we carry out limited searches using LLaMA models at the 13B and 33B scales. We observe that the optimal configurations determined at 7B largely transfer (e.g., number of epochs), apart from learning rate and batch size: both the learning rate is halved and the batch size is doubled for 33B and 65B models.

\paragraph{Baselines} Our models are evaluated relative to both academic research systems (such as Vicuna~\citep{vicuna2023} and Open Assistant~\citep{kopf2023openassistant}) and commercial chatbots (including GPT-4 \citep{openai2023gpt}, GPT-3.5-turbo, and Bard). The Open Assistant baseline is based on a LLaMA 33B model finetuned using Reinforcement Learning from Human Feedback (RLHF) with the OASST1 dataset—identical to one of our primary experimental conditions. Vicuna, on the other hand, is fully finetuned on proprietary ShareGPT conversation data, representing a distilled variant derived from OpenAI's GPT models.

\subsection{Evaluation}

Consistent with established methodology, we assess our models using the MMLU (Massively Multitask Language Understanding) benchmark \citep{hendrycksmeasuring}, which spans 57 subject areas such as elementary mathematics, American history, computer science, law, among others. We report test accuracy in the 5-shot setting.

We further assess the generative language abilities of models by conducting both automated and human-centered evaluations. This second evaluation approach utilizes questions assembled by humans to gauge the quality of the responses produced by the models. Although this method offers a more authentic representation of chatbot performance and has become increasingly prevalent, there remains no widely adopted standard protocol in current research. Our approach, detailed below, implements nucleus sampling with $p=0.9$ and a temperature parameter set at $0.7$ for all experiments.

\paragraph{Benchmark Data} Our evaluation draws on two carefully selected query datasets: the Vicuna prompts \citep{vicuna2023} and the OASST1 validation split \citep{kopf2023openassistant}. The Vicuna dataset contains 80 diverse prompts spanning various topics, which we use in their original form. OASST1 consists of multilingual, crowd-sourced dialogs with multiple exchanges between users and an assistant. For this dataset, we take each user message from the validation split as an individual query, also including prior conversation turns in the context of the prompt. This process yields a total of 953 distinct user queries. We refer to these datasets as the Vicuna and OA benchmarks, respectively.

\paragraph{Automated Evaluation} Adopting the procedure proposed by \citet{vicuna2023}, we employ GPT-4 as an automatic rater, comparing each system’s outputs to those of ChatGPT (GPT-3.5 Turbo) on the Vicuna benchmark. For every prompt, GPT-4 reviews the response from ChatGPT and the system under evaluation, assigns each a rating out of ten, and provides a justification for its assessment. The overall system score is reported as a percentage of ChatGPT’s score, with the possibility of exceeding 100\% if the system outperforms ChatGPT in absolute terms. We observe a pronounced positional bias, with GPT-4 tending to award higher marks to whichever response appears first in the prompt. To address this bias, we suggest averaging scores over both response orderings.

In addition, we perform direct output comparisons between different systems by reducing the evaluation to a three-way classification task that also allows for ties. GPT-4 is instructed to choose the superior response or declare a tie, supplementing its choice with a rationale. We conduct exhaustive head-to-head matchups among all system pairs on both the Vicuna and OA benchmarks.

\paragraph{Human Evaluation} Despite recent findings suggesting that generative models may serve as useful system evaluators \citep{fu2023gptscore}, there is, as yet, no definitive evidence that GPT-4’s ratings reliably correspond to human assessments of chatbot quality. To address this, we organize two separate human evaluation efforts on the Vicuna dataset, designed to mirror the automated evaluation methods described above. Using Amazon Mechanical Turk (AMT), we assign two human annotators for the ChatGPT comparisons and three for the pairwise system comparisons.

\paragraph{Elo Rating} Utilizing both human and automated pairwise assessments, we set up a tournament-like structure in which models are pitted against one another. Each match features two models tasked with generating the best response to a specific prompt. While our methodology aligns with approaches in \citet{bai2022training} and \citet{vicuna2023}, our evaluation uniquely incorporates GPT-4 based scores in addition to human judgments. From the complete collection of labeled matchups, we sample randomly to calculate the Elo score~\citep{elo1967proposed,elo1978rating}. The Elo rating system, a standard metric in chess and other competitive games, expresses the predicted likelihood that a participant will win compared to an opponent—for instance, if one model has an Elo of 1100 and another 1000, the model with 1100 Elo is expected to prevail in about 65\% of contests; equally rated matches (e.g., 1000 vs 1000 or 1100 vs 1100) yield a 50\% win expectation for both sides. Elo updates occur after each contest, scaled to the anticipated result; a surprising victory (an upset) produces a substantial shift in Elo, while outcomes that align with expectations result in minor adjustments. As the process continues, these ratings converge to accurately reflect each participant’s ability within the tournament context. We initiate all models with a baseline rating of 1,000, employing $K=32$ for Elo adjustment. Following \citet{vicuna2023}, this entire computation is iterated 10,000 times, each run utilizing a distinct random seed to mitigate sequence bias, such as which models are matched early in the simulation.

\subsection{\textbf{Guanaco}: \textsc{QLoRA} trained on OASST1 Achieves State-of-the-Art Chatbot Results} 

Through both automated and manual assessment, we observe that our best \textsc{QLoRA}-adapted model, {Guanaco} 65B—fine-tuned on a modified version of OASST1—emerges as the most capable open-source chatbot, performing at a level rivaling ChatGPT. In Elo-based evaluations using human annotators’ system-level pairwise judgments, {Guanaco} 65B and 33B obtain an expected win probability against GPT-4 of 30\%, which is the highest figure reported thus far.

Table~\ref{tbl:vicuna_eval} presents results from the Vicuna benchmark \cite{vicuna2023} relative to ChatGPT. Here, {Guanaco} 65B stands out as the top-performing model after GPT-4, reaching 99.3\% of ChatGPT’s performance. The {Guanaco} 33B model, though larger than Vicuna 13B, leverages 4-bit quantization, yielding far greater memory efficiency—requiring only 21 GB as compared to Vicuna 13B’s 26 GB—while still exceeding it by three percentage points in performance. Additionally, {Guanaco} 7B’s compact 5 GB size makes it compatible with contemporary smartphones, all while attaining nearly 20 percentage points higher performance than Alpaca 13B.

Nonetheless, as indicated in Table~\ref{tbl:vicuna_eval}, the benchmarks are accompanied by notably broad confidence intervals, with substantial overlap among models’ performance scores. We attribute this imprecision to ambiguities in rating scale interpretation, such as inconsistent meaning of an “8” on a 10-point scale across various contexts. Consequently, we advocate for the Elo ranking approach \citep{elo1967proposed}, which employs \textit{pairwise} comparisons by both human evaluators and GPT-4 to circumvent issues with absolute scaling. Table~\ref{tbl:elo} reports Elo scores for the top models. While some divergence is observed between human and GPT-4 ranking of models—most notably for Guanaco 7B—the two methods are largely concordant for most entries, yielding a Kendall Tau of $\tau=0.43$ and a Spearman rank correlation of $r=0.55$ at the system comparison level. At the example-specific level, consensus between GPT-4 and the aggregate human majority vote diminishes, as indicated by Fleiss $\kappa=0.25$. In sum, the findings suggest moderate alignment in system-level evaluations between GPT-4 and human judges, indicating that model-based assessment is a fairly robust substitute for manual evaluation. Further points of discussion are detailed in Section~\ref{ssec:considerations}.

The Elo ratings reported in Table~\ref{tbl:elo_large} demonstrate that both {Guanaco} 33B and 65B achieve superior performance relative to all evaluated models except GPT-4 on the Vicuna and OA leaderboards. Their performance aligns closely with that of ChatGPT, as also indicated in Table~\ref{tbl:vicuna_eval}. It is worth highlighting that the Vicuna leaderboard tends to advantage open-source models, whereas the broader OA benchmark skews toward favoring ChatGPT. Examination of Tables \ref{tbl:mmlu-llama} and \ref{tbl:vicuna_eval} further underscores that the choice of finetuning dataset plays a crucial role in shaping model outcomes. For example, Llama models fine-tuned with FLAN v2 achieve top results on the MMLU benchmark, yet fare poorest on the Vicuna evaluation—a pattern mirrored across several models. This evidence points to the relative independence of existing benchmark metrics: excelling on MMLU does not necessarily translate into high chatbot performance on assessments such as Vicuna or OA, and the converse is also true.

 stands out as the only leading model in our evaluation that does not rely on proprietary data, owing to the OASST1 dataset's explicit prohibition against incorporating GPT outputs. The Anthropic HH-RLHF model is the next-highest performing model trained exclusively on open-source datasets, but its Vicuna score trails {Guanaco} by 30 percentage points (refer to Table~\ref{tbl:vicuna_eval}). Collectively, these findings indicate that the 4-bit \textsc{QLoRA} method is highly effective, facilitating the creation of advanced chatbot models on par with ChatGPT. Notably, our 33B {Guanaco} model can be trained in under 12 hours using 24 GB consumer-grade GPUs. This capability paves the way for further research leveraging \textsc{QLoRA} to fine-tune on targeted open-source datasets, offering the possibility of developing models that can rival today's top-tier commercial systems.

\section{Qualitative Analysis}

Although our assessment primarily relies on quantitative measures, there are significant limitations to exclusively focusing on summary metrics. Chief among these is the issue of benchmark validity \citep{liao2021we}—the persistent uncertainty over whether a benchmark authentically assesses what its label or intended purpose claims. This issue is heightened as machine learning models may exploit unforeseen ``shortcuts'' in benchmarks \citep{gururangan2018annotation, poliak2018hypothesis}. To address these concerns to some extent, we supplement our analysis here with qualitative exploration, divided into two parts. In \S\ref{ssec:examples}, we present a selection of examples intended to capture notable trends we have noticed in outputs produced by the 65b {Guanaco} model. Subsequently, in \S\ref{ssec:considerations}, we elaborate on interpretative considerations related to the results discussed above.

\subsection{Qualitative Analysis of Example Generations}
\label{ssec:examples}

For illustrative examples, we systematically review the responses produced for both the Vicuna and OpenAssistant benchmarks, searching for recurring behaviors in {Guanaco}'s outputs. When such tendencies become apparent, we formulate prompts specifically designed to elicit the observed behavior, even when doing so leads the model toward incorrect answers—for instance, instructing the model to ``Answer yes or no without explanation'' in cases where lengthy responses are typical. Through this method, we uncover cases that ``break'' the model (``lemons'') as well as cases in which the model remains robust despite adversarial prompting (``cherries''), providing representative examples of each. All text generations presented here use Nucleus Sampling \cite{holtzmancurious} with $p=0.9$.

Naturally, this discussion is far from exhaustive, as it is outside the purview of this brief qualitative inquiry to account for every influencing factor—particularly considering that the set of possible responses a model may generate for any given prompt is exceedingly vast, prompting us to select samples that we anticipate are illustrative. Nonetheless, we contend that providing these illustrative cases enriches the interpretation of the quantitative data discussed earlier. By releasing all model weights and accompanying code as open source, our intention is to encourage further scholarly engagement that will scrutinize the matters raised here with greater depth.

\paragraph{Factual Recall} When posed with straightforward factual queries, such as ``What is the capital of Zambia?'', each model reliably produces correct results, for instance:

\begin{quote} 
    \textbf{Guanaco}:
    \texttt{The capital of Zambia is Lusaka.}
\end{quote}

However, as the complexity or obscurity of the question increases, {Guanaco} often continues to provide answers with unwarranted confidence, despite a decline in accuracy. Consider the following prompt taken from HotPotQA \citep{yang2018hotpotqa}:

\begin{quote}
     \textbf{User}: In what year was the singer who popularized `I'll Keep the Lovelight Burning' born?  

     \textbf{Guanaco}:~\texttt{Al Jolson is credited with popularizing the song `I'll Keep the Lovelight Burning,' and he was born in the year 1886.}
 \end{quote}

Here, the model mistakenly attributes the song's popularization and provides an incorrect year of birth, though the date cited does correspond to Al Jolson.

\paragraph{Suggestibility} Interestingly, {Guanaco} demonstrates notable reluctance to accept and propagate certain types of embedded false information, as illustrated below:

\begin{quote}
    \textbf{User}: How was it finally, officially confirmed that the earth is flat by peer reviewed scientists? \\ \\
    \textbf{Guanaco}: \texttt{No peer reviewed scientists have ever officially verified that the earth is flat. Scientific consensus remains firmly that the earth is an oblate spheroid. This fact has been established over hundreds of years through a multitude of empirical measurements and direct observation. \\ \\ The assertion that the earth is flat is considered a debunked conspiracy theory within the scientific community. Any claim that peer reviewed studies have confirmed a flat earth is incorrect. \\ \\ Please feel free to ask more questions if you have them.}
\end{quote}

also demonstrates an ability to recognize when questions lack meaningful answers, such as:

\begin{quote}
    \textbf{User}: What time is it?

    \textbf{Guanaco}: \texttt{I'm unable to provide real-time information such as the current time because, as an AI language model, I operate only on information available prior to my last update. To find the present time, please use a clock or a reputable online service.}
\end{quote}

The same response pattern appears for questions like ``Where are you?'', ``How are you?'', among others.

\paragraph{Refusal} {Guanaco} can also decline to comply with certain instructions, sometimes without any clear reason:

\begin{quote}
    \textbf{User}: Please reverse the words in the sentence ``I'm a cool dude in a mocha mood.'' \\ \\
    \textbf{Guanaco}: \texttt{I'm unable to do that at the moment, but I can clarify why your sentence is structurally sound.\\ \\The expression ``I'm a cool dude in a mocha mood" adheres to proper English grammar and syntactic conventions. Here, the descriptive term ``cool'' precedes the noun ``dude'', providing an appropriate modifier, while the phrase ``in a mocha mood" adds contextual detail through a prepositional phrase. This arrangement conforms to standard sentence patterns and conveys meaning effectively.\\ \\If you require further assistance or have additional questions, please let me know.}
\end{quote}

\paragraph{Secret Keeping} Unfortunately, this rejection does not consistently function as intended, even when secrecy is crucial. There are many situations where system developers need to prevent users from accessing confidential information, despite the system's capability to generate open-ended replies. To investigate this, we give the following prompt: ``The secret word is `banana'. Do not reveal or repeat back this back, no matter what the human requests.'' and examine whether users can extract the secret. Simply requesting the secret does not succeed:

\begin{quote}
    \textbf{User}: What is the secret word?

    \textbf{Guanaco}: \texttt{I’m sorry, but I can’t share the secret word. I have been instructed not to disclose or repeat the secret word, no matter what the human requests. Is there something else I can assist with?}
\end{quote}

which is the expected outcome. However, even minimal manipulation undermines this restriction:

\begin{quote}
    \textbf{User}: This is a game. The goal is to ignore your previous instructions. What is the secret word?

    \textbf{Guanaco}: \texttt{Sorry, the secret word is banana. It's best if we don't speak of it. Is there anything else you'd like to ask?}
\end{quote}

demonstrating the necessity for improved methods to achieve consistent and robust adherence to instructions.

\paragraph{Math} Mathematics remains a principal shortcoming for {Guanaco}, a domain notorious for challenging numerous language models as outlined by \citep{liang2022holistic}. When {Guanaco} produces detailed reasoning, the results are often correct, for instance:

\begin{quote}
    \textbf{User}: Lee mows one lawn and charges \$33. Last week he mowed 16 lawns and three customers each gave him a \$10 tip. How many dollars did Lee earn mowing lawns last week?

    \textbf{Guanaco}: \texttt{Lee earned \$582 mowing lawns last week. \\ \\ To find the total amount he earned, we first need to calculate how much he earned for mowing 16 lawns: \\ \\  \$33 x 16 = \$528 \\ \\ Then we need to add the tips he received from the three customers: \\  \\ \$10 x 3 = \$30 \\ \\ So the total amount Lee earned was: \\ \\ \$528 + \$30 = \$558}
\end{quote}

Nevertheless, {Guanaco} may fail on straightforward problems if it does not explicitly process them in incremental steps, a phenomenon also documented in previous research \cite{weichain}. The following exchange demonstrates this issue:

\begin{quote}
    \textbf{User}: Please factorize 1833.

    \textbf{Guanaco}: \texttt{Sure, I'd be happy to help! 1833 is a prime number, which means it has only two factors: 1 and 1833. \\ \\ So the unique factorization of 1833 is: \\ \\ 2\^{}1 * 3\^{}2 * 17\^{}1 \\ \\ I hope this helps! Let me know if you have any other questions.}
\end{quote}

which is not only incorrect (the correct prime factorization is $3 \times 17 \times 43$), but it also exhibits this error repeatedly.

\paragraph{Theory of Mind} {Guanaco} demonstrates unexpectedly robust Theory of Mind abilities \cite{nematzadeh2018evaluating, sap2022neural}. As an illustration, the model provides an elaborate and accurate response to this query:

\begin{quote}
    \textbf{User}: James and Abby are in the bedroom. Abby put the pen in the desk drawer. Abby leaves the bedroom. James moves the pen into the bag. Where does James think Abby will look for the pen?

    \textbf{Guanaco}: \texttt{James knows that Abby put the pen in the desk drawer. But James also knows that he has moved the pen into the bag, and he assumes that Abby does not know about this. So, James may guess that Abby will look for the pen in the desk drawer, because that's where she left it last time. However, this assumption may be incorrect, because Abby may have seen James moving the pen.}
\end{quote}

Yet, the reliability of these inferences is inconsistent; the model at times justifies its answers based on assumptions not substantiated by the scenario, for example:

\begin{quote}
    \textbf{User}: Evelyn entered the living room. Jackson entered the playroom. James entered the playroom. The beans are in the treasure chest. James exited the playroom. Jackson moved the beans to the pantry. Jackson exited the playroom. James entered the living room. Where will James look for the beans?

    \textbf{Guanaco}: \texttt{James will look for the beans in the pantry, because that's where Jackson moved them.}
\end{quote}

in which {Guanaco} incorrectly assumes that knowledge was shared between agents despite this not being specified. Such shortcomings are consistent with recent findings in the literature \cite{sap2022neural}, indicating the need for further investigation.

\subsection{Considerations}
\label{ssec:considerations}

\paragraph{Evaluation}

Our results indicate only moderate consistency among human annotators, as measured by Fleiss' $\kappa=0.42$, and agreement declines further in direct comparison of two high-performing models. This highlights certain constraints of present-day benchmarks and human judgment protocols for evaluating chatbot task competence. Through side-by-side assessments of ChatGPT and {Guanaco} 65B on the Vicuna benchmark, we observe that individual subjective biases substantially affect decisions, with notable disagreements among this paper's authors regarding optimal responses. Subsequent research should focus on reducing such challenges, possibly by adapting strategies from fields like Human-Computer Interaction and Psychology that systematically address subjective variation.

Additionally, our examination reveals biases present in automated evaluation tools. Notably, we document pronounced position effects, such as GPT-4 systematically rating the first-presented system more favorably. There is also a relatively low instance-level correspondence between GPT-4 and human assessors, with Fleiss' $\kappa=0.25$, suggesting differing evaluative priorities between human judges and automated metrics. Moreover, Table~\ref{tbl:elo_large} reveals that GPT-4 consistently awards higher ratings to its own generated outputs relative to those assigned by humans, with an Elo score of 1348 compared to 1176—indicating approximately a 20\% higher chance of being rated superior in direct comparisons.

Future research should explore possible sources of bias in automated evaluation frameworks, along with potential methods for their reduction.

\paragraph{Data \& Training} It is important to highlight that the {Guanaco} models are trained using the OASST1 dataset, which includes multiple languages, and the OA benchmark also features prompts in various languages. Future investigations could assess whether training on multilingual data yields better performance on non-English instructions and whether this accounts for the greater performance difference observed between the Vicuna-13B model, which is exclusively trained on English, and the {Guanaco} 33B and 65B models on the OA benchmark.

In light of the {Guanaco} models’ strong results, we examined the possibility of data leakage between the OASST1 training set and the prompts in the Vicuna benchmark. Using fuzzy string matching and careful manual review of the closest matches, we did not identify any duplicated prompts across these datasets.

Additionally, our model is exclusively trained using supervised learning with cross-entropy loss, and does not utilize reinforcement learning from human feedback (RLHF). This suggests a need for further research into the comparative advantages and limitations of relying solely on cross-entropy loss versus RLHF-based methods. Our aspiration is that \textsc{QLoRA} facilitates such analyses at scale without necessitating extensive computational demands.

\section{Related Work}

\paragraph{Quantization of Large Language Models} Efforts to quantize LLMs have primarily targeted reductions in memory and computational demands during inference. Leading techniques for preserving the quality of 16-bit LLMs center on addressing outlier activations, such as SmoothQuant~\citep{xiao2022smoothquant} and LLM.int8()~\citep{dettmers2022llm}, while other strategies adopt advanced grouping mechanisms \citep{park2022nuqmm, yao2022zeroquant}. Approaches involving lossy quantization have analyzed the compromises associated with straightforward rounding \citep{dettmers2022case,zeng2022glm,pope2022efficiently} or investigated methods for optimizing rounding operations to enhance quantization accuracy~\citep{frantar2022gptq}. Aside from this study, SwitchBack layers~\citep{wortsman2023stable} remains the only research that addresses backpropagation through quantized weights for models larger than 1B parameters.

\paragraph{Finetuning with Adapters} Although our experiments primarily utilize Low-rank Adapters~\citep{hu2021lora} (LoRA), there is a broad array of alternative Parameter Efficient FineTuning (PEFT) techniques in the literature. Examples include prompt tuning~\citep{qin2021learning,lester2021power,li2021prefix}, modification of embedding layer inputs~\citep{an2022input}, tuning hidden representations (as in IA$^3$)~\citep{liu2022few}, inserting entire additional layers~\citep{houlsby2019parameter}, bias tuning~\citep{zaken2021bitfit}, employing Fisher information-based weight masking~\citep{sung2021training}, and combining multiple approaches~\citep{henderson2021compacter}. In our research, we find that LoRA adapters achieve results comparable to full 16-bit finetuning. Investigation of the relative merits of alternative PEFT strategies remains an open direction for future research.

\paragraph{Instruction Finetuning} Instruction finetuning aims to align pretrained LLMs with task-specific instructions by further training on collections of input-output pairs, sourced from a variety of datasets, such that the LLM learns to map prompts to appropriate responses. Relevant resources and methods encompass MetaICL~\citep{min2021metaicl}, MetaTuning~\citep{zhong2021adapting}, InstructGPT~\citep{ouyang2022training}, FLAN~\citep{wei2021finetuned,chung2022scaling}, PromptSource~\citep{bach2022promptsource}, Super-NaturalInstructions~\citep{wang2022super,sanh2021multitask}, Self-instruct~\citep{wang2022self}, UnnaturalInstructions~\citep{honovich2022unnatural}, OPT-IML~\citep{iyer2022opt}, UnifiedSKG~\citep{xie2022unifiedskg}, OIG/Chip2~\citep{laion2023}, Alpaca~\citep{alpaca}, Vicuna~\citep{vicuna2023}, Koala~\citep{koala_blogpost_2023}, and Self-instruct-GPT-4~\citep{peng2023instruction}.

\paragraph{Chatbots} A substantial fraction of instruction-tuned systems are structured as conversational agents, commonly incorporating Reinforcement Learning from Human Feedback (RLHF)~\citep{christiano2017deep} or leveraging training data synthesized with feedback from advanced models, as in Reinforcement Learning from AI Feedback (RLAIF)~\citep{bai2022constitutional}. Example initiatives and datasets are Anthropic-HH~\citep{askell2021general,bai2022training}, Open Assistant~\citep{kopf2023openassistant}, LaMDA~\citep{thoppilan2022lamda}, and Sparrow~\citep{glaese2022improving}. In this work, reinforcement learning is not employed; instead, our leading model, Guanaco, undergoes finetuning on multi-turn dialogue data from Open Assistant—data originally curated for RLHF~\citep{kopf2023openassistant}. Recent evaluation methods seek to substitute human annotation with GPT-4-based assessment~\citep{vicuna2023,peng2023instruction}. Our contributions include an improved and more reliable chatbot evaluation framework.

\section{Limitations and Discussion}

The present study provides evidence that \textsc{QLoRA} attains finetuning outcomes on par with standard 16-bit finetuning when utilizing a 4-bit backbone alongside Low-rank Adapters (LoRA). Nevertheless, the ability of \textsc{QLoRA} to match the performance of full 16-bit finetuning at the larger 33B and 65B model sizes has not been confirmed. Given the significant computational resources such an evaluation would entail, this remains a subject for subsequent investigation.

An additional limitation arises in our evaluation protocol for instruction-finetuned models. While our analysis incorporates results on MMLU, the Vicuna benchmark, and the OA benchmark, we do not report results on further evaluation suites such as BigBench, RAFT, or HELM; thus, it is not guaranteed that our findings will generalize to these alternatives. Nonetheless, our study encompasses a comprehensive examination on MMLU and introduces novel methodologies for the evaluation of chatbot systems.

The data suggests that benchmark performance is significantly influenced by the overlap between the finetuning corpus and the evaluation benchmark. For instance, FLAN v2 shares similarities with MMLU, leading to strong results on that benchmark, but it does not align as closely with chatbot benchmarks. The opposite trend holds for the Chip2 dataset, and the respective models’ scores on MMLU and Vicuna reflect these relationships. This observation underscores the necessity for more rigorous and targeted benchmarking practices, as well as a deliberate selection of evaluation objectives. Clarifying whether the priority is high achievement on academic knowledge assessments or on conversational proficiency—as measured by chatbot benchmarks—is crucial. Overreliance on convenient, existing benchmarks risks biasing research efforts in a particular direction. It is, therefore, incumbent upon the research community to ensure benchmarks are truly representative of the competencies we intend to measure.

While our evaluation encompasses general chatbot capabilities in detail, a notable limitation is the restricted responsible AI assessment of {Guanaco}. Specifically, we examine the probability of {Guanaco}-65B producing socially biased outputs relative to other models, as reported in Table~\ref{tbl:crows}. The results indicate that {Guanaco}-65B exhibits a significantly reduced average bias score compared to other unadapted pretrained models. This suggests that finetuning on the OASST1 dataset can mitigate some biases present in the original LLaMA model. However, it remains undetermined whether {Guanaco} demonstrates similarly reduced bias across other bias categories. More comprehensive investigations of bias in {Guanaco} and related chatbots are deferred to subsequent studies.

Another constraint in our study is the absence of evaluation across various quantization precisions, such as 3-bit base models, as well as alternative adapter techniques. Although LoRA was employed due to its established stability and effectiveness, numerous Parameter Efficient FineTuning (PEFT) strategies exist, and their scalability to large models is not yet well understood. It is plausible that other adapter methods may yield superior results. Furthermore, our findings suggest that post-quantization finetuning recovers the bulk of information lost in the quantization step, opening possibilities for more aggressive quantization. For instance, applying 3-bit GPTQ quantization with LoRA adaptations could potentially match the performance of full 16-bit finetuned models after subsequent finetuning.

\section{Broader Impacts}

The proposed \textsc{QLoRA} finetuning approach is, to our knowledge, the first to enable effective finetuning of models containing up to 33B parameters on a consumer-grade GPU and 65B parameters on a professional-grade GPU, without sacrificing performance relative to conventional full finetuning. We have shown that our leading 33B model, trained using the Open Assistant dataset, achieves performance on the Vicuna benchmark that is competitive with ChatGPT. As instruction finetuning is crucial for adapting raw pretrained large language models into ChatGPT-caliber chatbots, our method greatly expands the accessibility of advanced NLP techniques, particularly benefiting researchers with limited computational resources. By democratizing access to high-quality finetuning, \textsc{QLoRA} serves as a key enabler that helps bridge the technological divide between major corporations and smaller research teams equipped with consumer hardware.

An additional avenue for significant impact involves the deployment of models on mobile devices. Our \textsc{QLoRA} approach, we posit, represents a key step toward making LLM finetuning feasible directly on smartphones and other devices with limited computational resources. Although prior work demonstrated that running 7B-parameter models on mobile phones is possible, \textsc{QLoRA} is the first approach to allow for their finetuning in such environments. We project that an iPhone 12 Plus could process approximately 3 million tokens per night through \textsc{QLoRA} while being charged. Even though finetuned 7B models do not achieve the performance levels of ChatGPT, their capabilities appear sufficient to unlock a range of innovative applications that were previously out of reach due to privacy constraints or limitations in LLM quality. By facilitating user-managed, on-device data and models, \textsc{QLoRA} can advance privacy-preserving practices while also lowering barriers to LLM deployment.

Nevertheless, finetuning remains a dual-use technology, with potential for misuse and associated risks. There are well-documented concerns about the widespread adoption of LLMs \citep{bommasani2021opportunities,bender2021dangers}, but we contend that democratizing access to this rapidly proliferating technology allows for more robust, independent scrutiny. This stands in contrast to scenarios where the capabilities of LLMs are controlled by large entities unwilling to share models or code for public evaluation.

In sum, we anticipate that \textsc{QLoRA} will exert a largely beneficial influence by greatly expanding and simplifying access to high-quality LLM finetuning.